

\section{Method}
\label{sec:method}

Given a few color images \eg $\{I_i\}_{i=1,2,3}$, where \( I_i \in \mathbb{R}^{H \times W \times 3} \), our goal is to obtain a 3D reconstruction of the observed scene or object efficiently. This reconstruction entails a 3D triangle mesh $\mathcal{S}$ and novel view synthesis capability. We achieve this objective by learning a generalizable feed-forward model \( \Phi \) that can map the input images to planar surface elements of any target view in a single forward pass (Fig.\ref{fig:pipe}). These $2$D primitives enable efficient depth and color volume rendering via perspective accurate Gaussian Splatting \cite{huang20242d}. An explicit mesh can be obtained from the splatted depths through the TSDF algorithm \cite{zhou2018open3d}, while color rendering enables real time novel view synthesis. 

%Given a few color images \eg \( I = \{I_i\}_{i=1,2,3} \), where \( I_i \in \mathbb{R}^{H \times W \times 3} \), our objective is to learn a mapping \( \Phi \) that outputs $2$D Gaussian parameters for both geometry and appearance. To achieve this, we utilize the pretrained MASt3R features which encode robust correspondence features to our MVS backbone to predict the necessary attributes for $2$D gaussians. First, we provide a brief overview of Gaussian Splatting based approaches in Section \ref{subsec:method_gs} and MASt3R in Section \ref{MASt3R}. We then detail in Section \ref{sec:feat_merge} how we utilize MASt3R features, obtained from image pairs as input to the feature extractor backbone of our MVS-based architecture which accepts any number of input images ($3$ for our sparse $3$D reconstruction setup) to predict $2$D Gaussian splatting parameters to predict geometrically accurate depth maps, followed by our losses in Section \ref{losses}.



\subsection{Background: Gaussian Splatting}
\label{subsec:method_gs}
\textbf{3D Gaussian Splatting (3DGS)}~\cite{kerbl20233d} represents a scene as a set of 3D anisotropic Gaussians \( \{ \mathcal{G} \} \), each defined by centroid \( \mathbf{x} \), scale \( \mathbf{S} \), rotation \( \mathbf{R} \), opacity \( \mathbf{\alpha} \), and SH-encoded color \( \mathbf{c} \). The covariance is:  
\begin{equation}
  \mathbf{\Sigma} = \mathbf{R} \mathbf{S} \mathbf{S}^T \mathbf{R}^T.
\end{equation} 
It can be projected to screen space via \( \mathbf{J} \) and \( \mathbf{W} \):  
\begin{equation}
  \mathbf{\Sigma'} = \mathbf{J} \mathbf{W} \mathbf{\Sigma} \mathbf{W}^T \mathbf{J}^T.
\end{equation}
Here, \( \mathbf{W} \) is the world-to-camera transformation matrix, and \( \mathbf{J} \) is the local affine transformation matrix that projects the Gaussian onto the image plane. Rendering projects pre-ordered Gaussians to the image plane and blends them via \( \alpha \)-compositing:  
\begin{equation}
\label{eq:alpha_blend}
 \mathbf{C} = \sum_{i \in \mathbf{N}} \mathbf{\alpha}_i' \mathbf{c}_i \prod_{j=1}^{i-1}(1-\mathbf{\alpha}_j'),
\end{equation} 
where the adjusted opacity is:  
\begin{equation}
\label{eq:alpha_mod}
\mathbf{\alpha}_i' = \mathbf{\alpha}_i e^{-\frac{1}{2}(\mathbf{y'}-\mathbf{x'})^T \mathbf{\Sigma'^{-1}}(\mathbf{y'}-\mathbf{x'})}.
\end{equation} 
Here, \(\mathbf{y'}\) is the pixel and \(\mathbf{x'}\) is the Gaussian splat centre's ($\mathbf{x}$) projection onto the 2D image plane.

3D primitives can induce ambiguity when viewed from different viewpoints. The 2D footprint can correspond to multiple possible 3D primitive configurations. Additionally, the perspective affine approximation used to project 3D Gaussian primitives onto the image plane becomes less accurate as points move away from the primitive center. These elements lead to 3D inconsistencies that 2D primitives help alleviating. 

\noindent \textbf{2D Gaussian Splatting (2DGS)}~\cite{huang20242d} represents primitives as planar 2D Gaussians, improving multi-view depth consistency while maintaining novel view performance. Each Gaussian in screen space is defined as:

\begin{equation}
\mathcal{G}(\boldsymbol{x}) = \exp\left(-\frac{u(\boldsymbol{x})^2 + v(\boldsymbol{x})^2}{2}\right),
\end{equation}

where $u(\boldsymbol{x})$ and $v(\boldsymbol{x})$ are the UV coordinates of the ray-splat intersection on the local tangent plane of the splat, and $\boldsymbol{x}$ is the screen space pixel coordinate. The ray-splat intersection has a closed form expression involving the homography matrix:
\begin{equation}
\boldsymbol{H} = \begin{bmatrix}
    s_u \boldsymbol{t}_u & s_v \boldsymbol{t}_v & \boldsymbol{0} & \boldsymbol{p} \\
    0 & 0 & 0 & 1
\end{bmatrix},
\end{equation}

where $s_u$, $s_v$, \ie $\boldsymbol{s} \in \mathbb{R}^{2}$ are learnable scaling factors, and $\boldsymbol{t}_u$, $\boldsymbol{t}_v$ are tangential vectors that can be derived from the learnable splat rotation $\boldsymbol{q} \in \mathbb{R}^{4}$, and $\boldsymbol{p \in \mathbb{R}^{3}}$ is the splat's learnable center in world space. Thereafter, $\mathbf{\alpha}_i' = \mathbf{\alpha}_i\mathcal{G}(\boldsymbol{x})$ as in Equ.\ref{eq:alpha_mod} with a learnable opacity parameter $\boldsymbol{\alpha} \in \mathbb{R}^{1}$ per splat. Alpha-blending is performed as in Equ.\ref{eq:alpha_blend} to get the rendered color from the depth-ordered primitives, using learnable colors $\boldsymbol{c}$ per splat. This approach improves multi-view depth consistency and 3D reconstruction quality by accurately modeling perspective in splat rendering, as opposed to the 3DGS based geometry modelling.


%\subsection{DUSt3R and MASt3R} \label{MASt3R}
%Our approach aims to utilize features of a strong prior in the form of the recently introduced $3$D foundation model, MASt3R~\cite{leroy2024grounding}, which is a new 3D-aware matching model built on DUSt3R and outputs dense local features that are robust to extreme viewpoint changes and significantly outperform current state-of-the-art methods. Trained on large datasets, it learns to establish correspondences in feature space between pairs of images of the same scene. This allows for multi-view stereo reconstruction by directly regressing 3D point clouds. Given images $\textbf{I}_1$ and $\textbf{I}_2$ of size $W \times H \times 3$, MASt3R outputs 3D points $\hat{X}_1$ and $\hat{X}_2$, all in the coordinate frame of the first image. It also outputs local dense feature descriptors $\hat{D}_1$ and $\hat{D}_2$, and confidence maps $C_1$ and $C_2$

% \begin{figure*}[t]
%     \centering
%     \includegraphics[width=1.0\linewidth]{figs/pipelinefig.pdf}
%     \vspace{-30pt}
%     \caption{Pipeline: Our approach takes input images, $I_{1}, I_{2}, I_{3}$,  and passes them through frozen MASt3R Feature Aggregation block to generate Aggregated features, $K_{1}, K_{2}, K_{3}$. The aggregated features are passed as input to MVS backbone to predict the 2D Gaussian splatting parameters. We employ predicted 2DGS parameters to perform sparse multiview reconstruction and novel view synthesis.}
%     \vspace{-15pt}
%     \label{fig:masterpipe}
% \end{figure*}

%The training objective $L_{pts}$ is:
%\begin{equation}
%    L_{pts} = \sum_{v \in \{1, 2\}} \sum_{i} C_{v}^{i} L_{regr}(v, i) - \gamma \log(C_{v}^{i}),
%\end{equation}
%where:
%\begin{equation}
%    L_{regr}(v, i) = \left\| \frac{1}{z} X_{v}^{i} - \frac{1}{\bar{z}} \hat{X}_{v}^{i} \right\|,
%\end{equation}
%Here, $C_{v}^{i}$ represents the confidence of depth estimation, helping to manage errors from ambiguous points. $\gamma$ is a hyperparameter controlling the weight of confidence, while $z$ and $\bar{z}$ are normalization factors, typically set to 1 for metric datasets. We use a pre-trained MASt3R model~\cite{leroy2024grounding} to extract our muti-view features. 

\subsection{Model}

Given $N$ input images $I_i$ and their camera poses $P_i$ (encompassing 3D rotation and translation), our model $\Phi$ predicts a pixel-aligned 2DGS Gaussian parameter map $O_t = \{\boldsymbol{s},\boldsymbol{q},\boldsymbol{p},\boldsymbol{\alpha},\boldsymbol{c}\} \in \mathbb{R}^{H\times W \times (2+4+3+1+3)}$ for a given target view $P_t$: 
\begin{equation}
O_t = \Phi(P_t|I_1,\ldots,I_N). 
\end{equation}
For test time reconstruction, TSDF depth fusion is performed on $N$  mean splatted depths. Each depth $D_i$ is splatted for view $P_i$ using the set of Gaussian primitives made of the inferred splat map $O_i=\Phi(P_i|I_1,\ldots,I_N)$. 

The Gaussian primitive attribute prediction pipeline follows the architecture in \cite{liu2024fast}, and we refer the reader to this latter for exhaustive details. An FPN network predicts a 2D feature map $f_i$ for each input image $I_i$. Homography based warping is used to warp feature maps to the target view $f_t^i$ using the relative poses $P_t$/$P_i$. Two parallel stages follow: a deepMVS branch and a pixel aligned prediction branch. The former follows classical deep depth from multi-view stereo architectures \cite{gu2020cascade}. Features $f_t^i$ are merged into a cost volume, which is processed into a probability volume with a 3D convolutional network, leading to a depth prediction. In the second branch, features $f_t^i$ are pooled into a target feature map $f_t$. A 2D convolutional network and an MLP map this feature map to Gaussian primitive attributes ${\{\boldsymbol{s},\boldsymbol{q},\boldsymbol{\alpha},\boldsymbol{c}\}}$. The unprojected deepMVS depth via camera intrinsics $K$ provides the 3D position attribute of the splats ${\{\boldsymbol{p}\}}$, thus forming the final full attribute map $O_t$. We also follow the hybrid rendering introduced in \cite{liu2024fast}.  

  




%Our architectural backbone, illustrated in Figure \ref{fig:masterpipe} is based on the MVSGaussian framework. It takes the $N$ feature maps $[K_{1}, K_{2}, K_{3}, ..., K_{N}]$ ($N = 3$ for our sparse reconstruction experiments) as input to predict the per-pixel $2$D Gaussian Splatting parameters for each image. To process these feature maps, we utilize a shared Feature Pyramid Network (FPN) that extracts multi-scale features.
%For each input feature map, the extracted features, together with the corresponding camera parameters $[K_i, R_i, t_i]_{i=1,2,3}$, are warped to a target viewpoint defined by the camera parameters $[K_t, R_t, t_t]$. The warped features are then used to construct a cost volume by computing their variance, which captures the consistency of multi-view features. This cost volume is subsequently passed through a $3$D convolutional network (CNN) for regularization, producing a probability volume that represents the depth likelihood at each pixel. The resulting depth probability distribution is used to weigh each depth hypothesis, allowing us to derive the final depth map. The feature maps from the FPN, together with the estimated depth map, are then warped into the target view. These warped features are aggregated and processed through a U-Net to predict the final $2$D Gaussian Splatting parameters, scale $s$, rotation $r$, color $c$, opacity $\alpha$, position $\mu$. 

\subsection{Additional Features}
MVSFormer++~\cite{cao2024mvsformer++} integrates 2D foundation model DINOv2~\cite{oquab2023dinov2} with cross-view attention and adaptive cost volume regularization to improve transformer-based Multi-View Stereo. In order to enrich the initial features $f_t^i$ of our pipeline, We follow this line to integrate and experiment with two types of deep visual features, namely monocular features from DINOv2~\cite{oquab2023dinov2} and multi-view features from MASt3R~\cite{leroy2024grounding}, and show that MASt3R's features enrich cost volume (which are 4D tensors that store matching costs across multiple views to estimate depth) features, thereby making the predicted depth, and consequently, the mesh reconstruction more accurate.
\vspace{-5pt}
\paragraph{Multi-view features}
\label{sec:feat_merge}
%MASt3R is a state-of-the-art 3D foundation model that solves 3d dowstream tasks such as point map prediction and correspondance generation using pairs on unposed images as input. It is made of a ViT encoder, a transformer decoder, in addition to task heads. We propose to use the features produced by the decoder. 
MASt3R is a state-of-the-art 3D foundation model trained on a large set of data, which learns to establish correspondences in feature space between pairs of images of the same scene. MASt3R takes two input images and predicts a dense pointmap and stereo feature image corresponding for each input image:
\begin{equation}
F^{j}_{i}, F^{i}_{j} = \mathcal{F}_{\text{MASt3R}}(I_{i}, I_{j}),
\end{equation}
where $F\in\mathbb{R}^{H\times W \times 24}$. 
We run master on all pairs made with our $N$ input images. Subsequently, we can define the feature of an image $I_i$ as the concatenation of the image with all its $N-1$ feature maps obtained in combination with the other views: $[I_i,F_i^k]_{k=1,k\neq i}^{k=N}$. For $N=3$, this writes  $[I_1,F_1^2,F_1^3]$. A simple way to inject this feature representation into our Gaussian attribute regression pipeline is by feeding it as input to the FPN network.

%In our problem setting, $N$ images of the same scene are given as input to predict the depth and scene parameters. We employ MASt3R to predict per-image feature maps between pairs formed from $N$ input images. %This serves as a strong generalizable prior to establish relations between pairs of images. 
%We generate $2\times N$ sets of features from $N$ input images using MASt3R which we use as input to our network along with the input images themselves. So, MASt3R, $M_{\theta}$,  takes a pair of input images, $I_{1}, I_{2}$ and predicts feature images, $F^{2}_{1}, F^{1}_{2}$,  as follows:            
%\begin{equation}
%\begin{split}
%    F^{2}_{1}, F^{1}_{2} &= M_{\theta}(I_{1}, I_{2}) \\
%    F^{3}_{1}, F^{1}_{3} &= M_{\theta}(I_{1}, I_{3}) \\
%    F^{3}_{2}, F^{2}_{3} &= M_{\theta}(I_{2}, I_{3})
%\end{split}
%\end{equation}


%For each input image, we concatenate feature images generated by pairing with other two images along with the image itself to generate input for our network. We concatenate feature images for each image along with input image to create the input for our approach. For an image, $I_{1}$, we concatenate its corresponding features from MASt3R to create the input, $K_{1}$, for our approach as follows:   
%\begin{equation}
%    K_{1}=[I_{1}, F^{2}_{1}, F^{3}_{1}]
%\end{equation}
%We extend this procedure for the other two images to generate the corresponding feature maps, denoted as $K_{2}$ and $K_{3}$ which we refer to as aggregated features. This procedure is illustrated in Figure \ref{fig:masteragg}.

%\begin{figure}
%    \centering
%    \includegraphics[width=1.0\linewidth]{figs/masteragg.pdf}
    % \vspace{-15pt}
%    \caption{MASt3R feature aggregation Block. We send three images into the MASt3R feature aggregation block to generate Aggregated features. We choose three pairs of two images to pass through frozen MASt3R pretrained model to get corresponding feature images, $F$. For each input image, the output feature images generated with other two images are concatenated with original input image to generate the aggregated features, $K$. We employ the aggregated features as input for our network. }
%    \vspace{-8pt}
%    \label{fig:masteragg}
%\end{figure}
\vspace{-5pt}
\paragraph{Monocular features} 
Deep foundational monocular features are more straightforward to inject in our 2DGS feed-forward regression pipeline. For DINOv2~\cite{oquab2023dinov2}, we extract the image feature map and up-sample it to the resolution of the input:
\begin{equation}
F_{i} = \mathcal{F}_{\text{DINO}}(I_{i}),
\end{equation}
where $F\in\mathbb{R}^{H\times W \times 384}$. The final feature input to the FPN network is the concatenation of the image and its monocular feature $[I_i,F_i]$.

%Then, as in \ref{sec:feat_merge}, we concatenate these features with the image itself to generate input for our network.




\subsection{Training Objective} 
\label{losses}

Our model $\Phi$ is trained in an end-to-end fashion using multi-view RGB images and ground truth depth as supervision, while deep feature networks ($\mathcal{F}$) are frozen for computational efficiency. We follow the multi-stage training procedure in \cite{liu2024fast}. A target image is differentiably reconstructed from $N$ neighboring input views. To optimize the generalizable framework, we employ a combination of the mean squared error (MSE) loss ($\mathcal{L}_{\mathrm{mse}}$), structural similarity index (SSIM) loss~\cite{ssim} ($\mathcal{L}_{\mathrm{ssim}}$), perceptual loss~\cite{lpips} between groundtruth image and rendered image ($\mathcal{L}_{\mathrm{perc}}$). Differently from \cite{liu2024fast}, we employ an $L_1$ depth loss ($\mathcal{L}_{\mathrm{depth}}$) between groundtruth depth and predicted depth. Also differently, we use a depth distortion ($\mathcal{L}_{d}$) and normal consistency ($\mathcal{L}_{n}$) losses to regularize our 2DGS output. The overall loss for each stage $k$ of the coarse-to-fine optimization framework described in \cite{liu2024fast} is formulated as:
\begin{equation}
\mathcal{L}^{k} = \mathcal{L}_{\mathrm{mse}} + \lambda_s \mathcal{L}_{\mathrm{ssim}} + \lambda_p \mathcal{L}_{\mathrm{perc}} + \lambda_\alpha \mathcal{L}_{d} + \lambda_\beta \mathcal{L}_{n} + \lambda_\gamma \mathcal{L}_{depth},
\label{eq:loss}
\end{equation}
where $\mathcal{L}_{d}$ and $\mathcal{L}_{n}$ are defined as:
\begin{equation}
\begin{split}
\mathcal{L}_{d} &= \sum_{i,j} \omega_i \omega_j |z_i - z_j|,\\
\mathcal{L}_{n} &= \sum_{i} \omega_i (1 - n_i^\mathrm{T} N).
\end{split}
\end{equation}
Here, $z_i$ and $z_j$ are the depths of splats along the viewing direction. $n_i$ is the surface normal of the splat, and $N$ is the normal from the rendered depth map. $\omega_i$ and $\omega_j$ are the blending weights. 
%Here, \todo{n, N, z, w, say what these things are briefly, make sure u sue the same syntax u tend to change symbols}. 
The depth distortion loss $\mathcal{L}_{d}$, as introduced initially in \cite{barron2021mip}, helps concentrating the weight distribution along the rays. The normal consistency loss, $\mathcal{L}_{n}$, ensures that the 2D splats are locally aligned with the actual surface geometry. Please refer to \cite{huang20242d} for more details regarding these losses.

Due to our MVS architecture being pyramidal in design, $\mathcal{L}^{k}$ represents the loss for the $k$-th stage, with $\lambda_s$, $\lambda_p$, $\lambda_\alpha$, $\lambda_\beta$, and $\lambda_\gamma$ denoting the respective weights of each loss term. The overall loss function is the sum of the losses from all stages, expressed as:
\begin{equation}
\label{eq:overall_loss}
\mathcal{L} = \sum_k \lambda^k \mathcal{L}^k,
\end{equation}
where $\lambda^k$ represents the weighting factor for the loss at stage $k$ (detailed in supplementary material). It is worth noting that unlike implicit methods which are unable to render full images and depths during training due to computation time constraints, we are able to apply the training losses on the entire images on account of the efficient Gaussian Splatting based rendering.